% Segmentation Pipeline section. Paste into your own thesis, or \input it.

\section{Network architecture}
\label{sec:pipeline-architecture}

CAS-Net \cite{dong2023casnet}, like every method in Table~\ref{tab:pipeline-benchmark} except
TotalSegmentator, is a 3D \textbf{U-Net} \cite{ronneberger2015unet}: an \textbf{encoder} that
repeatedly shrinks spatial resolution while growing feature depth, so deeper layers see more
context; a \textbf{decoder} that mirrors this back up to full resolution; and \textbf{skip
connections} that carry the encoder's fine spatial detail directly across to the matching decoder
level, so the output is both contextually informed and spatially precise.

CAS-Net is a five-level 3D U-Net with three added attention modules (Figure~\ref{fig:pipeline-casnet}),
aimed at a problem plain skip connections do not solve well: a coronary artery is thin and
branching, so an encoder feature and a decoder feature at the same location are not equally
reliable everywhere. Each module lets the network learn, per channel and per location, how much
to trust each signal (function names refer to \texttt{models/cas\_net.py}):

\begin{itemize}
  \item \textbf{SAFE}, at the bottleneck: re-processes the coarsest feature map with dilated
    convolutions at several scales, then self-attention, so distant voxels belonging to the same
    vessel can influence each other directly.
  \item \textbf{AGFF}, at every skip connection: computes a per-channel attention weight and uses
    it to scale the encoder's contribution before combining it with the decoder's, rather than
    concatenating both equally everywhere.
  \item \textbf{MSFA}, at the output: combines the decoder's output from all four resolution
    levels into the final prediction, rather than basing it on the finest level alone.
\end{itemize}

\begin{figure}[htbp]
  \centering
  \resizebox{\textwidth}{!}{%
  \begin{tikzpicture}[font=\small,
      enc/.style={draw, fill=blue!10, minimum width=15mm, minimum height=8mm},
      dec/.style={draw, fill=orange!10, minimum width=15mm, minimum height=8mm},
      mod/.style={draw, rounded corners, fill=green!10, minimum width=15mm, minimum height=7mm},
    ]
    \node[enc] (e0) at (0,0) {enc.\ input\\16 ch};
    \node[enc] (e1) at (0,-1.4) {enc.\ 1\\32 ch};
    \node[enc] (e2) at (0,-2.8) {enc.\ 2\\64 ch};
    \node[enc] (e3) at (0,-4.2) {enc.\ 3\\128 ch};
    \node[enc] (e4) at (0,-5.6) {enc.\ 4\\256 ch};
    \draw[-{Latex[length=2mm]}] (e0)--(e1); \draw[-{Latex[length=2mm]}] (e1)--(e2);
    \draw[-{Latex[length=2mm]}] (e2)--(e3); \draw[-{Latex[length=2mm]}] (e3)--(e4);

    \node[mod, fill=purple!12] (safe) at (2.6,-5.6) {SAFE\\(bottleneck\\attention)};
    \draw[-{Latex[length=2mm]}] (e4) -- (safe);

    \node[mod] (agff4) at (5.2,-4.2) {AGFF 4};
    \node[dec] (d4) at (7.8,-4.2) {dec.\ 4\\128 ch};
    \node[mod] (agff3) at (5.2,-2.8) {AGFF 3};
    \node[dec] (d3) at (7.8,-2.8) {dec.\ 3\\64 ch};
    \node[mod] (agff2) at (5.2,-1.4) {AGFF 2};
    \node[dec] (d2) at (7.8,-1.4) {dec.\ 2\\32 ch};
    \node[mod] (agff1) at (5.2,0) {AGFF 1};
    \node[dec] (d1) at (7.8,0) {dec.\ 1\\16 ch};

    \draw[-{Latex[length=2mm]}] (safe) -- (agff4); \draw[-{Latex[length=2mm]}] (agff4)--(d4);
    \draw[-{Latex[length=2mm]}] (d4) -- (agff3); \draw[-{Latex[length=2mm]}] (agff3)--(d3);
    \draw[-{Latex[length=2mm]}] (d3) -- (agff2); \draw[-{Latex[length=2mm]}] (agff2)--(d2);
    \draw[-{Latex[length=2mm]}] (d2) -- (agff1); \draw[-{Latex[length=2mm]}] (agff1)--(d1);

    \draw[-{Latex[length=2mm]}, dashed, gray!60!black] (e3) to[bend right=10] (agff4);
    \draw[-{Latex[length=2mm]}, dashed, gray!60!black] (e2) to[bend right=10] (agff3);
    \draw[-{Latex[length=2mm]}, dashed, gray!60!black] (e1) to[bend right=10] (agff2);
    \draw[-{Latex[length=2mm]}, dashed, gray!60!black] (e0) to[bend right=10] (agff1);

    \node[mod, fill=red!10, minimum height=16mm, align=center] (msfa) at (10.6,-2.1)
      {MSFA\\(fuses all\\4 decoder\\levels)};
    \draw[-{Latex[length=2mm]}] (d4) -- (msfa); \draw[-{Latex[length=2mm]}] (d3) -- (msfa);
    \draw[-{Latex[length=2mm]}] (d2) -- (msfa); \draw[-{Latex[length=2mm]}] (d1) -- (msfa);
    \node[draw, minimum width=15mm, minimum height=8mm, right=8mm of msfa] (out)
      {2-channel\\logits};
    \draw[-{Latex[length=2mm]}] (msfa) -- (out);
  \end{tikzpicture}%
  }
  \caption{CAS-Net \cite{dong2023casnet}: a five-level 3D U-Net (blue/orange, channel counts
    per level) with three attention modules added -- SAFE at the bottleneck, AGFF at every
    skip connection, and MSFA fusing all four decoder resolutions into the final prediction.}
  \label{fig:pipeline-casnet}
\end{figure}

The network's output is a two-channel volume of raw logits (background, lumen); nothing in the
model decides a hard yes/no answer -- that is deferred to postprocessing
(Section~\ref{sec:pipeline-inference}).
